%%%%%%%%%%%%%%%%%%%%%%%%%%%%%%%%%%%%%%%%%%%%%%%%%%%%%%%%%%%%%%%%%%%%%%%%%%%%%%%%%%
\begin{frame}[fragile]\frametitle{}
\begin{center}
{\Large K-Means Clustering: Customer Segmentation}

{\tiny (Ref: scikit-learn Clustering User Guide)}
\end{center}
\end{frame}

%%%%%%%%%%%%%%%%%%%%%%%%%%%%%%%%%%%%%%%%%%%%%%%%%%%%%%%%%%
\begin{frame}[fragile]\frametitle{Problem Info}
\begin{itemize}
\item Scenario: A retail company wants to segment its customers for targeted marketing
\item Features: Annual Income (k USD) and Spending Score (1--100)
\item Data: Synthetic dataset mimicking a retail loyalty program
\item Goal: Discover natural customer groups using K-Means; choose optimal K with the Elbow method
\end{itemize}
\end{frame}

%%%%%%%%%%%%%%%%%%%%%%%%%%%%%%%%%%%%%%%%%%%%%%%%%%%%%%%%%%
\begin{frame}[fragile]\frametitle{Import Libraries and Generate Data}
\begin{lstlisting}
import numpy as np
import pandas as pd
import matplotlib.pyplot as plt
from sklearn.cluster import KMeans
from sklearn.preprocessing import StandardScaler

np.random.seed(42)
# Simulate 3 customer segments: budget, standard, premium
centers = [(30, 20), (55, 55), (80, 85)]
X = np.vstack([
    np.random.randn(120, 2) * [8, 10] + c
    for c in centers])
df = pd.DataFrame(X, columns=['AnnualIncome', 'SpendingScore'])
print(df.describe())
\end{lstlisting}
\end{frame}

%%%%%%%%%%%%%%%%%%%%%%%%%%%%%%%%%%%%%%%%%%%%%%%%%%%%%%%%%%
\begin{frame}[fragile]\frametitle{Elbow Method: Choosing K}
\begin{itemize}
\item Inertia (WCSS, within-cluster sum of squares): the total squared distance of the points to their cluster centers
\item Look for the elbow: the K after which extra clusters help little
\end{itemize}
\begin{lstlisting}
scaler = StandardScaler()
X_scaled = scaler.fit_transform(df)

inertia = []
K_range = range(1, 10)
for k in K_range:
    km = KMeans(n_clusters=k, random_state=42, n_init=10)
    km.fit(X_scaled)
    inertia.append(km.inertia_)

plt.plot(K_range, inertia, marker='o')
plt.xlabel("Number of Clusters K")
plt.ylabel("Inertia (WCSS)")
plt.title("Elbow Method")
plt.grid(True); plt.show()
\end{lstlisting}
\end{frame}

%%%%%%%%%%%%%%%%%%%%%%%%%%%%%%%%%%%%%%%%%%%%%%%%%%%%%%%%%%
\begin{frame}[fragile]\frametitle{Fit K-Means with Optimal K}
\begin{lstlisting}
km = KMeans(n_clusters=3, random_state=42, n_init=10)
df['Segment'] = km.fit_predict(X_scaled)

# KMeans cluster numbers are arbitrary, so map them to
# Budget/Standard/Premium by sorting on mean income
order = df.groupby('Segment')['AnnualIncome'].mean().sort_values().index
labels = dict(zip(order, ['Budget', 'Standard', 'Premium']))
colors = {'Budget': '#e07b54', 'Standard': '#5b9bd5', 'Premium': '#70ad47'}

plt.figure(figsize=(7, 5))
for seg, name in labels.items():
    mask = df['Segment'] == seg
    plt.scatter(df.loc[mask, 'AnnualIncome'],
                df.loc[mask, 'SpendingScore'],
                label=name, color=colors[name], alpha=0.7)
plt.xlabel("Annual Income (k USD)")
plt.ylabel("Spending Score")
plt.title("Customer Segments (K-Means, K=3)")
plt.legend(); plt.tight_layout(); plt.show()
\end{lstlisting}
\end{frame}

%%%%%%%%%%%%%%%%%%%%%%%%%%%%%%%%%%%%%%%%%%%%%%%%%%%%%%%%%%
\begin{frame}[fragile]\frametitle{Fit K-Means with Optimal K: Intuition}
\begin{block}{Intuition}
KMeans cluster numbers (0, 1, 2, \ldots) are arbitrary and can come out in any order;
they carry no meaning by themselves. To turn them into business labels like ``Budget'' or
``Premium'', sort the clusters by a real statistic (here, mean income) after fitting,
rather than assuming cluster 0 is always the first segment.
\end{block}
\end{frame}

%%%%%%%%%%%%%%%%%%%%%%%%%%%%%%%%%%%%%%%%%%%%%%%%%%%%%%%%%%
\begin{frame}[fragile]\frametitle{Segment Profiles}
\begin{itemize}
\item \textbf{Budget}: low income, low spending; price-sensitive customers
\item \textbf{Standard}: mid income, mid spending; core customer base
\item \textbf{Premium}: high income, high spending; loyalty program targets
\end{itemize}
\begin{lstlisting}
profile = df.groupby('Segment')[
    ['AnnualIncome','SpendingScore']].mean()
profile.index = [labels[i] for i in profile.index]
print(profile.sort_values('AnnualIncome').round(1))
\end{lstlisting}
\end{frame}

%%%%%%%%%%%%%%%%%%%%%%%%%%%%%%%%%%%%%%%%%%%%%%%%%%%
\begin{frame}[fragile]\frametitle{Results: Elbow and Segments}
\begin{itemize}
\item Inertia drops sharply up to $K = 3$ (720, 217, 88) and then only slowly: the elbow is at $K = 3$
\item The three segments match the three groups that generated the data
\end{itemize}
\begin{lstlisting}
# inertia for K = 1, 2, 3, ...
720.0  216.9  87.7  76.5  65.2  54.8  48.2  41.4  37.3

# segment profiles: mean income, mean spending score
          AnnualIncome  SpendingScore
Budget            29.3           20.7
Standard          55.4           54.7
Premium           79.0           84.6
\end{lstlisting}
\end{frame}

%%%%%%%%%%%%%%%%%%%%%%%%%%%%%%%%%%%%%%%%%%%%%%%%%%%
\begin{frame}[fragile]\frametitle{Quick Check: Reading the Elbow}
\begin{block}{Think About It}
The inertia for $K = 1, 2, 3, 4, 5$ was 720, 217, 88, 77, 65. Which choice of $K$ is best supported, and why?
\begin{enumerate}
  \item[A.] $K = 9$, because it has the lowest inertia
  \item[B.] $K = 3$, because the big drops stop there and extra clusters help only a little
  \item[C.] $K = 1$, because it keeps all customers together
  \item[D.] Any $K$, because inertia does not depend on the number of clusters
\end{enumerate}
\end{block}
\end{frame}

%%%%%%%%%%%%%%%%%%%%%%%%%%%%%%%%%%%%%%%%%%%%%%%%%%%
\begin{frame}[fragile]\frametitle{Quick Check: Reading the Elbow}
\begin{block}{Answer}
\textbf{Correct: B.} Inertia always falls as $K$ grows and reaches zero when every point is its own cluster, so the lowest value (A) is a trap. Look for the elbow: the drop from 217 to 88 is large, and the next drops are small. Here that matches the three groups that generated the data. $K = 1$ (C) hides all structure, and inertia clearly depends on $K$ (D).
\end{block}
\end{frame}